\section{结果来源与统计口径核对}
本附录保留截至 2026 年 10 月 15 日的开发、历史复查与旧演示结果，并补充保存 10 月 16 日五算法紧凑演示及受限语义后端的详细验证。
完整开发、三批历史复查、筛选演示与网关应用分别报告，
不把不同标注范围、选择过程与用途合并为一个“总体成绩”。
本附录采用实现包含口径；网关接口资产在第四章单独说明。
以下后端对照、回归与旧筛选实验均为截至 2026 年 10 月 15 日的留存结果；
10 月 16 日混合演示复扫与后端验证详情在本附录末尾单列，第四章集中报告两项算法实验及两个前端应用。
报告构建程序重新从逐文件行计算 TP/FP/FN/TN，
并与原汇总核对，通过后生成表格和图像。
原实验报告、生成时间、源文件哈希及绘图脚本见附录 A。

\section{完整开发集与历史验证复查}
\subsection{总体结果}
\begin{table}[htbp]\centering\small
\caption{主集合实现包含结果；历史集合是已观察数据的复查}\label{tab:overall-results}
% 表格直接保存在本章，编译不依赖实验数据目录。
\resizebox{\linewidth}{!}{%
\begin{tabular}{@{}lrrlrr@{}}\toprule
集合 & 文件 & 正标签 & TP/FP/FN & 召回率 & 精确率\\\midrule
完整开发集 & 2344 & 1572 & 1486/0/86 & 94.53\% & 100.00\%\\
历史主验证集 & 2855 & 1455 & 684/6/771 & 47.01\% & 99.13\%\\
最近作者集 & 232 & 176 & 97/6/79 & 55.11\% & 94.17\%\\
RF 历史留出集 & 252 & 233 & 142/3/91 & 60.94\% & 97.93\%\\
\bottomrule\end{tabular}}
\end{table}

完整开发集检出 1486 个正标签，漏报 86 个，召回率为 94.53\%，
漏报率为 5.47\%，已标注范围内 FP 为 0。
这是已知家族开发回归结果，不表示未知实现达到同样水平。
历史主集检出 684 个、漏报 771 个，召回率为 47.01\%，漏报率为 52.99\%。
最近作者集 TP/FP/FN 为 97/6/79，RF 历史留出集为 142/3/91，
漏报率分别为 44.89\% 和 39.06\%。
名称中的 RF 来源于该集合的历史随机森林研究用途，
本表仅使用当前默认规则，未把 RF 输出混入预测。

开发与三个历史集合的标签总体准确率分别为 99.66\%、96.94\%、97.18\%、97.13\%。
这些数值包含大量已标注阴性标签；即使总体准确率超过 90\%，
历史主集仍漏掉一半以上的真实正标签。
因此“准确率达到 90\%”不能替代“召回率达到 90\%”或“漏报低于 10\%”的结论。
对真实未知实现，当前工具仍需改进覆盖，不能仅用准确率宣布目标达成。

开发状态为 2338 completed / 6 partial；
历史主集为 2821 / 34，最近作者集为 230 / 2，RF 留出集为 252 / 0。
所有部分完成输入保留在分母中；未把失败通道、预算限制或未知结构过滤出结果。

\subsection{算法级结果}
\begin{table}[htbp]\centering\small
\caption{十三种算法的开发与历史主集结果}\label{tab:algorithm-results}
% 表格直接保存在本章，编译不依赖实验数据目录。
\resizebox{\linewidth}{!}{%
\begin{tabular}{@{}llrlr@{}}\toprule
算法 & 开发 TP/FP/FN & 开发召回 & 历史 TP/FP/FN & 历史召回\\\midrule
AES & 142/0/3 & 97.93\% & 112/0/28 & 80.00\%\\
DES & 119/0/12 & 90.84\% & 37/0/38 & 49.33\%\\
RC4 & 100/0/8 & 92.59\% & 30/0/45 & 40.00\%\\
ChaCha20 & 136/0/4 & 97.14\% & 88/0/44 & 66.67\%\\
RSA & 61/0/3 & 95.31\% & 19/0/83 & 18.63\%\\
ECDSA & 59/0/3 & 95.16\% & 0/0/59 & 0.00\%\\
ECDH & 96/0/5 & 95.05\% & 14/0/101 & 12.17\%\\
MD5 & 108/0/0 & 100.00\% & 95/0/23 & 80.51\%\\
SHA-1 & 112/0/10 & 91.80\% & 95/6/59 & 61.69\%\\
SHA-256 & 155/0/14 & 91.72\% & 71/0/104 & 40.57\%\\
SM2 & 149/0/8 & 94.90\% & 0/0/114 & 0.00\%\\
SM3 & 130/0/15 & 89.66\% & 71/0/25 & 73.96\%\\
SM4 & 119/0/1 & 99.17\% & 52/0/48 & 52.00\%\\
\bottomrule\end{tabular}}
\end{table}

% 配图 P05：十三算法横向分组条形图，开发与历史使用不同颜色，共用 0--100% 轴。
\begin{figure}[htbp]\centering
\includegraphics[width=\linewidth]{05-algorithm-recall.png}
\caption{同一算法在开发与历史主集上的实现召回率。图中差距反映实现与构建覆盖差异，不是时间序列增益。}
\label{fig:algorithm-recall}
\end{figure}
% 配图注释（P05）：放结果统计图，不放流程图。建议横向分组条形图，算法名保持完整，
% 开发与历史使用青色和暖灰色；标出百分比，并在图注注明实现包含口径。

开发集上，MD5 的 108 个正标签全部检出；SM4 为 119/120，
但 SM3 仍漏 15/145，说明开发汇总良好也不代表每算法均低于 10\% 漏报。
历史主集的公钥缺口更明显：RSA 漏 83/102，ECDH 漏 101/115，
ECDSA 漏 59/59，SM2 漏 114/114。
这些算法使用共享大数、点运算和跨函数协议，
当前规则难以从陌生的无符号实现恢复足够完整的关系。
SHA-256 与 SM4 也存在优化表示和不同家族覆盖不足的问题。

历史 6 个 FP 位于 SHA-1 类别。共享 IV 或相近结构的区分仍需要改进，
不能因精确率整体较高而忽略其对具体资产清单的影响。
本章根据留存报告解释错误分布，未用历史预测重新调参后再回写同一份成绩。

\clearpage
\subsection{源码与程序载体的差异}
\begin{table}[htbp]\centering\small
\caption{按载体分组的实现包含结果；源码 API 检出不在本表混算}\label{tab:kind-results}
% 表格直接保存在本章，编译不依赖实验数据目录。
\resizebox{\linewidth}{!}{%
\begin{tabular}{@{}llrrrr@{}}\toprule
集合 & 类型 & 文件 & TP & FN & 召回率\\\midrule
完整开发集 & 程序 & 2190 & 1408 & 59 & 95.98\%\\
完整开发集 & 源码 & 154 & 78 & 27 & 74.29\%\\
历史主验证集 & 程序 & 2677 & 671 & 693 & 49.19\%\\
历史主验证集 & 源码 & 178 & 13 & 78 & 14.29\%\\
最近作者集 & 程序 & 210 & 91 & 68 & 57.23\%\\
最近作者集 & 源码 & 22 & 6 & 11 & 35.29\%\\
RF 历史留出集 & 程序 & 230 & 138 & 82 & 62.73\%\\
RF 历史留出集 & 源码 & 22 & 4 & 9 & 30.77\%\\
\bottomrule\end{tabular}}
\end{table}
% 配图 P06：左侧载体对比，右侧历史 TP/FN 堆叠条形，不使用饼图掩盖分母。
\begin{figure}[htbp]\centering
\includegraphics[width=\linewidth]{06-error-distribution.png}
\caption{载体层次的实现召回率与历史主集的算法正标签分布。不同集合的源码上下文与标签组成不同。}
\label{fig:error-distribution}
\end{figure}
% 配图注释（P06）：适合放载体对比和 TP/FN 统计图。必须将源码的“实现包含”与源码的“API 引用”区分，
% 不要用源码 API 演示效果解释所有源码实现的准确率。

源码并不在任何口径下都比二进制容易。
明确 API 调用可以直接定位，而历史源码实现可能缺少头文件、工程上下文或完整算法阶段；
这些输入的实现级真值与源码引用的任务不同。
对二进制，符号剥离、优化、不同端序、查表变体和跨函数数据依赖影响覆盖。
因此不能仅按文件类型删除漏报对象，再把剩余成绩解释为全量能力提高。

\section{Miasm 语义后端的收益与成本}
\subsection{接入前后的预测对照}
\begin{table}[htbp]\centering\small
\caption{Miasm 接入前后历史复查；三批逐文件预测均未改变}\label{tab:miasm-results}
\resizebox{\linewidth}{!}{%
\begin{tabular}{@{}lrrrl@{}}\toprule
集合 & 文件 & 接入前召回 & 接入后召回 & Miasm 新增标签\\\midrule
历史主集 & 2855 & 47.01\% & 47.01\% & 0\\
最近作者集 & 232 & 55.11\% & 55.11\% & 0\\
RF 历史留出集 & 252 & 60.94\% & 60.94\% & 0\\\bottomrule
\end{tabular}}\end{table}

% 配图 P07：接入前后成对柱形与候选状态计数，真实显示没有召回增益。
\begin{figure}[htbp]\centering
\includegraphics[width=\linewidth]{07-miasm-comparison.png}
\caption{Miasm 接入前后对照及候选状态分布。候选完成是片段层次的状态，不表示完整函数分析成功。}
\label{fig:miasm-results}
\end{figure}
% 配图注释（P07）：放成对结果柱形图与预算状态统计，前后柱必须保持真实高度。
% 不宜使用“语义增强后显著提升”的示意图，因为本轮实际新增算法标签为零。

开发对照从 TP1484/FN88 变化到 TP1486/FN86，新增 RC4、DES 各一标签，
原报告确认来自旧有 Unicorn 探针，不归因给 Miasm。
后端的价值目前是提供可审查的表达式和后续调用摘要基础，
尚未通过这些历史数据证明新增泛化收益。

\subsection{预算、原语与失败状态}
三批历史候选合计 713 completed、4487 partial、528 budget\_exhausted、5 failed。
5 个失败为 Thumb 上游指令语义错误，被限制在独立子进程内，没有产生算法发现。
合计输出通用原语或指令语义角色 1127 条；
这些记录仍标为原语，未被当成 1127 个密码函数或正标签。
阶段耗时之和为 5736.39 秒，对 3117 个可执行文件约为 1.84 秒/文件，
包含候选选择与子进程启动；其余 222 个输入为源码。

32 进程历史主集墙钟为 1885.91 秒，最近作者集为 176.07 秒，RF 留出集为 158.73 秒。
阶段耗时之和与批次墙钟是不同量，不能混用。
旧轮并发与调度条件并非完全一致，因此不把全部墙钟差异归因给一个后端。
当前结果支持保留后端的研究和诊断作用，也支持在小型现场演示中关闭它以降低成本。

\section{五算法、两架构筛选演示复扫}
\begin{table}[htbp]\centering\small
\caption{按已观察表现选择的演示子集；不是新的独立验证}\label{tab:selected-results}
% 表格直接保存在本章，编译不依赖实验数据目录。
\resizebox{\linewidth}{!}{%
\begin{tabular}{@{}llrrlr@{}}\toprule
集合 & 架构 & 文件 & 正标签 & TP/FP/FN & 漏报率\\\midrule
开发演示 & x86-64 & 91 & 43 & 43/0/0 & 0.00\%\\
开发演示 & ARM64 & 90 & 44 & 44/0/0 & 0.00\%\\
历史演示 & x86-64 & 131 & 45 & 42/0/3 & 6.67\%\\
历史演示 & ARM64 & 125 & 39 & 35/0/4 & 10.26\%\\
\bottomrule\end{tabular}}
\end{table}
开发侧 181 文件、87 个正标签，TP/FP/FN 为 87/0/0；
历史侧 256 文件、84 个正标签，为 77/0/7，合并漏报率 8.33\%。
历史 x86-64 为 42/0/3、漏报 6.67\%；ARM64 为 35/0/4、漏报 10.26\%。
所以“合并低于 10\%”不能转述为两个架构分别均低于 10\%。

% 配图 P11：筛选集合的分架构结果柱形，10% 虚线，只说明该演示配置。
\begin{figure}[htbp]\centering
\includegraphics[width=.9\linewidth]{11-selected-architecture.png}
\caption{筛选演示输入的分架构漏报率。对所选变体的复扫不消除先前按表现选择产生的偏差。}
\label{fig:selected-architecture}
\end{figure}
% 配图注释（P11）：放四组分架构柱形图，显示 0\%、0\%、6.67\%、10.26\%，
% 加 10\% 参考线并写明“筛选演示、非独立验证”；不可把 ARM64 的柱画到线下。

437 份输入复制后全部新扫，与选择阶段逐文件预测一致；
未删除所有变体均漏的家族，7 个历史漏报仍保留。
两批各 16 进程同时运行，墙钟分别为 105.76 和 199.91 秒，
不能相加称为整轮墙钟，也不等于单线程速度。
阴性对照分别为 94、172 份，阴性数量不会降低按正标签计算的 FNR。
总体工具代码和原始样本未因此删除，筛选配置也未用于替代前述完整历史结果。

\section{回归状态、局限与目标对照}
主引擎最终回归记录为 28 模块、240 个测试方法，27 模块通过，
\nolinkurl{test_function_signatures} 模块失败。
失败涉及 ARM 阴性样例的完成状态断言，留存对照指向已有 Unicorn 子进程异常，
算法预测仍为空。没有通过修改旧断言、扩大总时限或删除样例取得全绿。
Miasm 的 14 项针对性测试通过不能替代这一全回归失败状态。

\begin{table}[htbp]\centering\small
\caption{目标与本轮能支持的结论}\label{tab:goal-review}
\begin{tabularx}{\linewidth}{p{3.1cm}X}\toprule
目标或展示项 & 实际结论\\\midrule
不少于五种密码算法 & 具备 13 种核心身份和相关通道；逐算法召回率差异明显。\\
不少于三类输入载体 & 源码、程序、固件入口已实现；不等于所有格式与算法均充分覆盖。\\
不少于五类真实电力终端 & 截至 10 月 15 日未做五类验证；10 月 16 日新增五类自制固件的软件场景见第四章，仍未做实体设备或真实厂商验证。\\
90\% 识别目标 & 开发召回 94.53\%；历史主召回 47.01\%。不能仅以标签准确率宣布未知实现目标达成。\\
演示漏报低于 10\% & 筛选历史子集合并为 8.33\%；ARM64 单独为 10.26\%，且存在表现选择偏差。\\
量子迁移辅助 & 已实现公钥资产关注和人工待办；未实现完整 PQC 协议或自动代码替换。\\\bottomrule
\end{tabularx}\end{table}

有效性威胁主要来自已见家族开发、编译变体相关、
源码上下文边界、逐行标注范围和语义预算；
当前没有足够证据声称对任意未知设备、算法变体或加密镜像具备通用能力。
后续评测应在开发阶段完善 ABI 与跨函数摘要后重新冻结，
使用未见作者、固定配置和完整失败分母检验真实收益。

\section{五算法紧凑混合演示的补充结果}


五算法紧凑集合共 241 份输入，包括开发来源二进制 107 份、历史来源二进制 104 份和开发源码 30 份，包含 196 个五算法范围内正标签及 45 个范围内阴性输入。该集合用于核对固定演示材料的识别和报告行为，结果见表~\ref{tab:showcase-origins}。

\begin{table}[!htbp]\centering\small
\caption{紧凑演示集合按来源分组的实现识别结果}\label{tab:showcase-origins}
% 表格直接保存在本章，编译不依赖实验数据目录。
\resizebox{\linewidth}{!}{%
\begin{tabular}{@{}lrrlrr@{}}\toprule
来源 & 文件 & 正标签 & TP/FP/FN & 召回率 & 精确率\\\midrule
开发二进制 & 107 & 87 & 87/0/0 & 100.00\% & 100.00\%\\
历史二进制 & 104 & 84 & 77/0/7 & 91.67\% & 100.00\%\\
开发源码 & 30 & 25 & 25/0/0 & 100.00\% & 100.00\%\\
混合演示合计 & 241 & 196 & 189/0/7 & 96.43\% & 100.00\%\\
\bottomrule\end{tabular}}
\end{table}

综合 TP/FP/FN 为 189/0/7，召回率为 96.43\%，精确率为 100.00\%，漏报率为 3.57\%，F1 值为 98.18\%。包含已审核阴性标签的文件--算法标签准确率为 99.26\%。由于该指标受到阴性标签数量影响，主要结论仍以正标签计数、召回率和漏报率为依据。图~\ref{fig:showcase-results} 给出结果与来源组成。

% 配图 P12：保留实际结果卡片与来源构成图，画面标明筛选混合演示。
% 显示开发与历史来源、正标签及漏报数量，不将该图标题改为独立验证准确率。
\begin{figure}[!htbp]\centering
\includegraphics[width=\linewidth]{12-showcase-results.png}
\caption{五算法紧凑演示集合的来源构成与识别结果。}
\label{fig:showcase-results}
\end{figure}

表~\ref{tab:showcase-algorithms} 与表~\ref{tab:showcase-kinds} 分别给出算法和载体分组。七个漏报全部位于历史来源二进制中，涉及 AES、MD5、SHA-1 和 SM4；SM3 在该集合内全部检出。x86-64 与 ARM64 的召回率分别为 96.59\% 和 95.18\%，开发源码的 25 个正标签全部检出。该源码部分包含依据已有表现选择的成功样例，其结果不能用于解释所有源码实现的识别难度。

\begin{table}[!htbp]\centering\small
\caption{紧凑演示集合的逐算法实现识别结果}\label{tab:showcase-algorithms}
% 表格直接保存在本章，编译不依赖实验数据目录。
\resizebox{\linewidth}{!}{%
\begin{tabular}{@{}lrlrr@{}}\toprule
算法 & 正标签 & TP/FP/FN & 召回率 & 精确率\\\midrule
AES & 42 & 40/0/2 & 95.24\% & 100.00\%\\
MD5 & 37 & 36/0/1 & 97.30\% & 100.00\%\\
SHA-1 & 40 & 38/0/2 & 95.00\% & 100.00\%\\
SM3 & 40 & 40/0/0 & 100.00\% & 100.00\%\\
SM4 & 37 & 35/0/2 & 94.59\% & 100.00\%\\
\bottomrule\end{tabular}}
\end{table}

\begin{table}[!htbp]\centering\small
\caption{紧凑演示集合按载体与架构分组的结果}\label{tab:showcase-kinds}
% 表格直接保存在本章，编译不依赖实验数据目录。
\resizebox{\linewidth}{!}{%
\begin{tabular}{@{}lrrlr@{}}\toprule
载体 / 架构 & 文件 & 正标签 & TP/FP/FN & 召回率\\\midrule
开发源码 & 30 & 25 & 25/0/0 & 100.00\%\\
x86-64 二进制 & 108 & 88 & 85/0/3 & 96.59\%\\
ARM64 二进制 & 103 & 83 & 79/0/4 & 95.18\%\\
\bottomrule\end{tabular}}
\end{table}

% 配图 P13：保留五算法 TP/FN 堆叠条与三类载体召回率图。
% 使用实际复扫数据，保留历史二进制漏报和各组分母，不使用截断坐标轴放大差异。
\begin{figure}[!htbp]\centering
\includegraphics[width=\linewidth]{13-showcase-coverage.png}
\caption{紧凑演示集合按算法与载体分组的识别结果。}
\label{fig:showcase-coverage}
\end{figure}

本集合沿用原二进制精简材料中的全部 171 份五算法阳性输入，包含已有七个漏报；阴性对照由 266 份缩减为 40 份，再加入 25 份已有成功源码和 5 份源码阴性输入。清单在复扫前固定，未按新结果再次筛选。241 份输入均完成扫描，16 进程墙钟时间为 138.84 秒。该实验的较高覆盖同时受到算法范围、编译变体和样例选择影响，因此只作为紧凑展示的补充结果，不与十三算法回归合并评价。


\section{受限语义后端的详细验证结果}

\label{sec:semantic-backend-results}

受限语义后端分别从局部行为与文件级识别进行验证。局部检查覆盖 18 组指令夹具和 208 组数值对照，并包含直接分支、路径合流、循环收敛、内存别名与预算退出等情形。已保存检查均通过，表明后端在这些片段和输入条件下能够保留预期的字宽与依赖关系。相同分支表达式可以在合流后继续使用，冲突状态则转为未知；当控制流、调用或预算条件不足时，后端保留限制信息，不将临时表达式视为完整函数结果。

在局部正反例中，摘要能够向既有运算识别逻辑提供 ChaCha 四分之一轮结构所需的运算关系。该结果只涉及算法局部结构，不能替代完整轮数、状态组织和调用条件的判断。进位、借位和宽乘等大数原语同样只作为后续推理材料，不直接赋予 RSA、ECDH、ECDSA 或 SM2 身份。

文件级工程对照沿用保存的三批历史复查，共 3339 份输入，其中可执行文件 3117 份、源码 222 份。该对照保留原历史架构范围，独立于第四章限定双架构的已知实现覆盖实验。接入前后，三批逐文件算法预测均未改变，Miasm 直接新增算法标签为零。开发对照中另外增加的 RC4 和 DES 标签，经证据核查来自原有行为探针，未归于 Miasm。文件级结果的候选状态与成本汇总见表~\ref{tab:semantic-backend-summary}；各批算法计数见附录表~\ref{tab:miasm-results}。

\begin{table}[!htbp]
  \centering
  \small
  \caption{受限语义后端在三批历史复查中的候选状态与分析开销}
  \label{tab:semantic-backend-summary}
  \begin{tabularx}{\linewidth}{@{}p{4.6cm}X@{}}
    \toprule
    观测项目 & 结果\\
    \midrule
    可执行文件数 & 3117 份\\
    候选完成 / 部分完成 & 713 / 4487\\
    候选预算耗尽 / 失败 & 528 / 5\\
    通用原语与指令语义记录 & 1127 条\\
    Miasm 直接新增算法标签 & 0\\
    后端阶段耗时之和 & 5736.39 秒\\
    平均每可执行文件阶段开销 & 约 1.84 秒，含候选选择与进程启动\\
    \bottomrule
  \end{tabularx}
\end{table}

候选状态对应局部分析任务，同一文件可以产生多个候选，不能将其数量解释为完整函数或文件数。五个失败均为 Thumb 上游指令语义错误，已限制在独立进程内，未形成算法发现。1127 条原语记录为进一步构造调用和算法摘要提供了可审查材料，但原语数量本身不构成算法召回收益。

由此可见，有限控制流、保守状态合并和表达式摘要已具备局部可验证性，当前接入方式尚未在留存文件级对照中增加算法检出，并引入额外分析成本。该结论支持将后端保留为可选的研究与诊断功能。后续收益仍需通过明确调用约定、输出缓冲区和协议关系的摘要及新的配对实验确认。上述阶段耗时之和不等于并行批次墙钟时间，也不能据其推导整套工具的平均响应时延。
